Was sind die Verantwortlichkeiten des Vorstandes in Bezug auf das ERM

Gesetz zur Kontrolle und Transparenz von Unternehmen §91 AktienG. 

"Der Vorstand hat geeignete Maßnahmen zu treffen, insbesondere ein
Überwachungssystem einzurichten, damit den Fortbestand der Gesell-
schaft gefährdende Entwicklungen früh erkannt werden."


Zudem §289 besagt, dass das Risikomanagementsystem im Lagebericht der Gesellschaft eingeführt wird, die nach §317 HGB auch Gegenstand der Abschlussprüfung sind.
Als Prüfungsstandard hat sich der IDW PS 340 etabliert.

Insbesondere hat der Vorstand dann folgende Tätigkeiten nachzugehen

- Die Entwicklung einer angemessenen Risikokultur, die im Unternehmen gelebt
und fortlaufend weiterentwickelt wird,
- die Verabschiedung von angemessenen unternehmensindividuellen Wesent-
lichkeitsgrenzen zur Bestimmung aller wesentlichen Risiken anhand geeigne-
ter und nachvollziehbarer Kriterien dem Risikoprofil


Was ist die Risikofunktion und Three lines of Defense 

Die Risikofunktion dient zur Umsetzung des Risikomanagements und besteht in kleineren Unternehmen meist aus einer in größeren auch aus mehreren Personen

- Beratung der Geschäftsleitung hinsichtlich Risiko,
- Untersuchung des Gesamtrisikos des Unternehmens unter Beachtung der Ab-
hängigkeiten der Einzelrisiken untereinander
-Erster Anlaufpunkt zur Meldung neuer Risiken 


Im Konzept der drei Verteidigungslinien (Three Lines of Defense) werden abgegrenzt:
- Risikonehmende Einheiten tragen die Verantwortung für die Steuerung der ei-
genen Profitabilität und des eigenen Risikoprofils (hinsichtlich Einzelrisiken).
- Zentrale risikoüberwachende Einheiten überwachen und unterstützen die risi-
konehmenden Einheiten und berichten die Risikolage an die Geschäftsleitung.
- Die Interne Revision prüft unabhängig das Risikomanagementsystem des Un-
ternehmens inklusive der beiden anderen Verteidigungslinien.


Zur der zweiten Linie gehören in Versicherungsunternehmen auch die VmF. und Compliance.


Was sind die Mitglieder des ERM-Ausschuses

Der Ausschuss dient zur Interaktion der Einheiten, welche auf Grund der Funktionstrennung/Unabhängigkeit notwendig ist.

- Vorstand 
- Risikofunktion 
- Vertreter operativer Einheiten

Welche Rolle spielt die interne Revision im ERM

Die Interne Revision prüft die Funktionsfähigkeit des Risikomanagementsystems
und unterscheidet dabei Aufbauprüfung und Funktionsprüfung.

Die Aufbauprüfung besteht aus
- Methodik,
- Organisation,
- Anpassungsfähigkeit,
- Integration in bestehende Überwachungs- und Führungssysteme.

Im Rahmen der Funktionsprüfung wird untersucht, ob das System im betrachteten
abgelaufenen Zeitraum tatsächlich wie vorgegeben konkret umgesetzt wurde und
die vorgesehenen Resultate erzielt wurden.

Konsistenz der Risikostrategie mit der Gesamtstrategie der Organisation,
- Konkretisierung in Bezug auf die Ableitung operativer Risikosteuerungsmaß-
nahmen,

Hinsichtlich der Risikoanalyse und -bewertung führt die IR Stichprobenartig Kontrollen durch 
bei qualitativen Analysen werden die Annahmen geprüft, bei quantitativen die unterliegenden Daten.

Hinsichtlich der Risikosteuerung wird geprüft 

- Beschreibung und Implementierung der definierten Indikatoren, Steuerungsmaßnahmen, Kontrollen und Überwachungsmaßnahmen
- Wirtschaftlichkeit gewählter Maßnahmen und Kontrollen zur Steuerung der
identifizierten Risiken

Hinsichtlich der Kommunikation und Berichterstattung:
- Festgelegte Rahmenbedingungen für die Berichterstattung. 
Was (Limits, Indikatoren) in welchen Kategorien wird wem (Adressat) wie (Medium), wann (Frequenz) berichtet.

- Für Regel- und Ad-hoc-Berichterstattung müssen der Kommunikationsprozess,
die jeweiligen Verantwortlichen für die Berichterstattung sowie die Berichts-
empfänger bestimmt sein. Ebenso für Gesetzliche Meldepflicht 


Was sind wichtige Bestandteile der Unternehmenskultur für Risikomanagement 

- Ein Führungsstil, der die Mitarbeiter konsultativ einbezieht.
- Übernahme von Verantwortung statt Schuldzuweisung.
- Eine gute interne Kommunikation/Wissensaustausch


Ein häufig anzutreffendes Phänomen in Unternehmen ist, dass (über)optimistische
Grundhaltungen zu unrealistischen Einschätzungen führen. Beispielsweise können
in Projektplanungen Risiken ausgeblendet und Chancen überwertet werden.

Ein zentrales Element der Risikokultur ist ein ausgewogenes Vergütungssystem.
Das Eingehen unangemessener Risiken darf nicht incentiviert werden. Tantiemen
sollten daher am langfristigen Unternehmenserfolg ausgerichtet werd




Definieren Sie das Risikokapital, seine Spezialisierung und Bewertungsmethoden 


Das in einem Unternehmen vorhandene Risikokapital sind die Mittel, die dem Un-
ternehmen zur Verfügung stehen, um Verluste aus unerwarteten Entwicklungen zu
tragen. 

Risikokapital = Vermögen - Verpflichtungen 

Welches Vermögen/Verpflichtungen genutzt werden darf hängt vom Standard ab (z.B. IRS, Solvency II etc.)

Das ökonomisch notwendige Risikokapital ist das Kapital, das erforderlich ist,
um bei ökonomischer Bewertung mögliche Verluste bei einer gegebenen Risikoto-
leranz über einen gegebenen Zeitraum zu bedecken.

Das betrieblich notwendige Risikokapital ist eine Spezialform des
ökonomisch notwendigen Risikokapitals. Hierbei beschreibt die Verlustfunktion den ökonomischen Wertverlust eines Unternehmens, also die Veränderung der Differenz
aus ökonomischem Wert der Vermögensgegenstände und ökonomischem Wert der Verpflichtungen.

Ratingkapital ist ein von Ratingagenturen ermitteltes benötigtes Risikokapital. Da
gute Ratingnoten eine große Finanzstärke signalisieren, wird Ratingkapital zu einem
sehr hohen Sicherheitsniveau ermittelt. Sie sind jedoch wenig detailliert und unterliegen höherer Fehlerquoten.

Solvenzkapital ist das regulatorisch vorgeschriebene Kapital, das ein Versiche-
rungsunternehmen vorhalten muss, um sein Geschäft betreiben zu dürfen.

Standardbewertungsansätze der Vermögenswerte für das Risikokapital sind
- run-off basis (Rückstellungen erlauben bei Auflösung des Konzerns vollständige Bestandsabwicklung)
- going-Concern basis (ähnlich wie oben, allerdings unter Annahme dass der Konzern weiterhin besteht)
- Referenzunternehmensbasis (Rückstellungen reichen aus damit ein drittes Vers. den Bestand damit abwickeln könnte)

Das Kapital wird zudem klassifiziert:
Methode der Kapitalschichtung:

- Rückstellungenskapital
- Risikokapital 
- Exzess Kapital 

Sowie dessen Ursprung 
- Aktionäre/Investoren 
- Echtes Fremdkapital (selten)
- Versicherungsnehmer zuzurechnendes Fremdkapital (Überschussfonds)

Dazu existiert noch das Tier System von Solvency II.



Was war die Motivation von Solvency II und was sind grob die drei Säulen 

1999 hat die EU-Komission bestimmt den bis dahin gültige regel-basierte Ansatz zur Ermittlung des benötigten Solvenzkapitals sollte durch einen
prinzipienbasierten Ansatz zu ersetzen.
Dies sollte indirekt folgende Ziele erreichen:
- EU-weit harmonisierten Aufsichts-
systems, das Wettbewerbsneutralität sicherstellt, die Entwicklung effizienter
Versicherungsmärkte fördert und einen ausreichenden Schutz der Interessen
der Versicherungsnehmer gewährleistet.

- Risikosensitive Eigenmittelanforderungen in Abhängigkeit vom Risikoprofil des Unternehmens 
- Anreize zur Entwicklung von internen Modellen
- Stärkung der Kontrollmechanismen des Marktes durch transparente und um-
fassende Berichtserstattung

Säule I (quantitative Anforderungen)
- Eigenkapitalanforderungen 
- Tier-System 
- Markt-konsistente Bewertung
Säule II (qualitative Anforderungen)
- Grundsätzen und Methoden des aufsichtsbehördlichen Überprüfungsverfah-
rens (SRP)
- ERM/Controlling, Governance
Säule III (Offenlegungs- und Berichts-
pflichten)


Die Richtlinien/Verordnungen der EU 289-291 AEUV, sind dabei die Quellen.

Auf nationaler Ebene sind das Versicherungsaufsichtsgesetz (VAG) und die Ausle-
gungsentscheidungen und Rundschreiben (MaGo) der BaFin die zentralen rechtli-
chen Grundlagen.



Erläutern sie detaillierter die Säulen von Solvency II


Säule I
Vermögenswerte und Verbindlichkeiten sind mit dem beizulegenden Zeitwert
(Fair Value/Market Value) anzusetzen. 

Vert. Rückstellungen haben keinen liquiden Referenzmarkt => Best Estimate + Risk margin

- Solvency Capital Requirement (SCR) 
- V@R mit Level 0.995 Zeitperiode 1 Jahr
- Minimum Capital Requirement (MCR)
- V@R mit Level 0.85 Zeitperiode 1 Jahr
Auf Basis eines internen Model oder der Bilanz wird über die Standard Formel 


Säule II:
Was nicht quantitativ abgedeckt werden kann, in einem wirksamen Risikomanagementsystem, fällt unter Säule II.

- Eine angemessene und transparente Aufbau- und Ablauforganisation mit klar
definierten Zuständigkeiten,
- eine angemessene Funktionstrennung
- klare Berichtslinien

Deshalb definiert sie die Schlüsselfunktionen 
- Risikomanagement (unabhängige Risikocontrollingfunktion)
- Compliance
- versicherungsmathematische Funktion
- Interne Revision 

Zudem bestimmt sie die
unternehmenseigene Risiko- und Solvabilitätsbeurteilung /
Own Risk and Solvency Assessment (ORSA). 
Über den ORSA der Säule II ist ein Bericht bei der Aufsicht einzureichen.

Säule III:
Ziel der auffsichtsrechtlichen Veröffentlichungs- und Berichtspflichten
sind die Stärkung der Marktmechanismen und der Marktdisziplin, sowie eine hinreicheinde Informationsbasis für die Aufsichtsbehörden.
Vorgeschrieben sind insbesondere
Solvency and Financial Conditions Report (SFCR):
-  Geschäftsstruktur und – strategie
-  finanziellen Situation und Ertragslage
-  Bewertungsansätzen für Kapital  & Verbindlichkeiten
-  SCR
-  Risikomanagement
-  ggf.  Angaben zum internen  Model
ergänzend ist der Regular Supervisory Report (RSR)mit weiteren Details für die Aufsichtsbehörde.

Hinzukommen rein quantiative Zusammenfassungen, sogenannte Templates.



Geben Sie eine Übersicht von Basel III und dessen Beziehung zu Solvency II 

Die Baseler Eigenkapitalempfehlung (Basel II) ist 2006 entstanden.
Wesentliches Ziel der Eigenkapitalregelung für Kreditinstitute ist es, 
die Kapitalanforderungen an Banken stärker am eingegangenen Risiko auszurichten.
Sie besteht ebenfalls aus einem 3 Säulen Modell.

Die wesentlichen Risiken die in Säule I abgedeckt werden sind
- Kreditrisiko,
- Marktrisiko 
- operationelles Risiko

Die Mindestkapitalanforderung für das Kreditrisiko wird mit einem 
faktorbasierten Ansatz als Faktor*Summe der risikogewichteten Aktiva bestimmt
Die Bewertung kann auf externen Ratings oder Internal Rating Based Approach (IRBA)
erfolgen.

Das Marktrisiko wird ebenfalls mit einem faktorbasierten Ansatz bestimmt. 

Zur Ermittlung der Mindestkapitalanforderung werden die drei Kapitalanforderungen addiert.

Säule II beschreibt qualitative Regelugen für das Risikomanagement
Anstelle des ORSA von Solvency II tritt dabei das 
Interne Kapitaladäquanzverfahren (Internal Capital Adequacy Assessment Process ICAAP).

Basel III ist eine Reform der Basel II motiviert durch die Finanzkrise 2007
Wesentliche Erneuerungen sind
- Stärkung von Qualität und Quantität des Eigenkapitals
- Einführung Verschuldungsobergrenze 
- 2 neue Liquiditätskennziffern sowie (LCR, NSFR)
- Verbesserung der Corporate Governance und Transparenz.

Im Detail sind dies 
- erhöhte Anforderungen an das Kernkapital und dessen Unterlegung, 
und zusätzliche Kapitalanforderungen auch durch einen antizyklischen Puffer
- Die Liquidity Coverage Ratio (LCR) soll sicherstellen, dass der Liquiditätsbedarf für die nächsten 30 Tage
gedeckt ist. 
- Die Net Stable Funding Ratio (NSFR) soll den langfristigen strukturellen Liquiditätsbedarf sicherstellen.

Ein wichtiger Unterschied zwischen Basel III und Solvency II ist dass die Risikobewertung bei ersterem
auf Basis der Rechnungslegung (Bilanz etc.) und die Risiken ohne Korellation, additiv behandelt werden.
Zweiteres nutzt eine extra erstellte Bilanz (SCFR Bestandteil).



Was ist die ebAV und dessen Beziehung zu Solvency II 

Das Kürzel steht für die Einrichtungen der betrieblichen Altersversorgung (EbAV)
gemäß EU Richtlinie 2016/234.

Betriebliche Altersversorgung (bAV) ist die Zusage von Leistungen der 
Alters-, Invaliditäts- oder Hinterbliebenenversorgung 
durch einen Arbeitgeber zugunsten seiner Arbeitnehmer.

Rechtsgrundlage ist in DE das Betriebsrentengesetz (BetrAVG).
Nicht alle Formen der bAV fallen unter Solvency II und werden deshalb gesondert behandelt. 
Insbesondere sind das die Pensionskassen und Pensionsfonds.

Bei Pensionskassen handelt es sich um Lebensversicherer, die wegfallendes Erwerbseinkommen versichern
Pensionsfonds erbringen Leistungen im Wege des Kapitaldeckungsverfahrens, wobei die Höhe der Leistungen und Beiträge
nicht immer fest zugesagt werden dürfen (im Gegensatz zur klassischen Versicherungsverträgen, Parallelen zur ges. Rentenver.)



Beschreiben sie die analogen Regelungen (zu Solvency II) der ebAV 

Säule I:
Es ist keine Solvabilitätsübersicht auf Marktwertbasis aufzustellen.
Bewertung der versicherungstechnischen Rückstellungen folgt dem Vorsichtsprinzip.
Es gilt insbesondere bezüglich der zulässigen Höchstrechnungszinsen und den verwendeten biometrischen Rechnungsgrundlagen.

Die anrechnungsfähige Eigenmittel werden in der ebAV, Solvabilitätsspanne genannt und sind:  

X Proz. der Deckungsrückstellung + Y Proz. des riskierten Kapitals

X = 
4, Standard,
1,  fondsgebundenen Versicherungen fuer ueber 5 Jahre
25, der Verwaltungskosten, im Falle von unter 5 Jahre

Y = 
0.3, Standard
0.15, 3 - 5 Jahre (Todesfallversicherungen)
0.1, 1 - 2 Jahre (Todesfallversicherungen)

die zurechenbaren Eigenmittel sind in jedem Fall bei 
- Aktiengesellschaften, eingezahlte Grundkapital - Betrag eigene Aktien,
- VVaG, eingezahlte Gründungsstock, Kapitalrücklagen und Gewinnrücklagen, ein Gewinnvortrag beziehungs-weise Verlustvortrag 
nach Abzug der auszuschüttenden Dividenden.
Zudem sind weitere Anrechnungen moeglich sofern genehmigt.

Von den sich so ergebenden Eigenmittel ist der Wert der immateriellen Vermögensgegenstände abzuziehen.

Säule II:


Das Governance System ist der Größenordnung, der Art, dem Umfang und der Komplexität der Tätigkeiten der EbAV 
und umfasst eine angemessene und transparente Organisationsstruktur mit einer klaren Zuweisung und 
angemessenen Trennung der Zuständigkeiten und ein wirksames System zur Gewährleistung der Übermittlung von Informationen.

Zudem sind Anlagevermögenswerte nach ökologischen, sozialen und Governance Faktoren nach Maßgabe der UN (United Nations) zu wählen.
Wie diese Faktoren in die Anlageentscheidungen einfließen ist offenzulegen werden, jedoch kann von diesen abgesehen werden
wenn Sie nicht im Verhältnis zur Größenordnung der Unternehmung stehen.

Die folgenden Schlüsselfunktionen müssen existieren:
- VmF
- Risikomanagement
- interne Revision

Die eigene Risikobeurteilung (ORSA) ist unter Berücksichtigung des Proportionali-
tätsgrundsatzes im Wesentlichen qualitativ durchzuführen und zu dokumentieren.
Adressat der Risikobeurteilung ist die Aufsichtsbehörde und findet alle 3 Jahre statt,

Säule 3:
Wesentliche Informationspflichten nach ebAV sind:
Jahresabschluss und einen jährlichen Lagebericht für jedes Versorgungssystem. 

Versorgungsanwärtern: 
- allgemeine Informationen zum Altersversorgungssystem 
- eine Leistungs-/Renteninformation 
- auf Anfrage bevor Rentenantritt Auskunft über die Auszahlungsoptionen erteilt werden.
Leistungsempfänger:
ähnlich wie Oben
- Leistungskürzungen müssen drei Monate vor Umsetzung.
Potenziellen Versorgungsanwärter:
- alle Optionen des Vorsorgesystems

Unabängig des Addressaten müssen bei Anlagerisiko und Anlageentscheidungen
Bericht über die Performance dieser innerhalb der letzen 5 Jahre geliefert werden.



Was ist das Proportionalitätsprinzip, Aspekte des Outsourcing und Notfallmanagement in Corporate Governance

Bei der Umsetzung der aufsichtsrechtlichen Anforderungen gilt das Proportionalitätsprinzip. 
Das bedeutet, dass der Wesensart, dem Umfang und der Komplexität der mit der Tätigkeit des Unternehmens 
einhergehenden Risiken Rechnung getragen werden muss. (Wesentlichkeit/Materialität)


Outsourcing muss so gestaltet sein, dass
- die Qualität des Governance-Systems nicht beeinträchtigt wird,
- das operationelle Risiko nicht übermäßig gesteigert wird,
- die Überwachung durch die Aufsicht nicht beeinträchtigt wird und
- dass keine Beeinträchtigungen für Versicherungsnehmer entstehen.

Das Notfallmanagement soll in Krisensituationen die Fortführung der Geschäftstätigkeit zu gewährleisten. 
Deshalb sind Notfallpläne zu führen für kritische Prozesse, welche die Geschäftstätigkeit gefährden können.
Diese sind regelmäßig zu testen.
